\documentclass[10pt,a4paper]{amsart}
\usepackage[margin=19mm]{geometry}
\usepackage{amsmath,amssymb,mathtools}
\usepackage[scr=rsfs,cal=euler]{mathalfa}
\usepackage{xfrac}
\usepackage[unicode,hidelinks,bookmarks=false]{hyperref}
\newcommand{\ie}{i.e.\ }
\newcommand{\cf}{cf.\ }
\newcommand{\resp}{respectively\ }
\newcommand{\Subsection}{\subsection}
\providecommand{\nobreakdash}{\nobreak\hbox{-}}
\setlength{\emergencystretch}{2em}
\multlinegap0pt
\usepackage[shortlabels]{enumitem}
\usepackage{mathtools}
\newtheorem{lemma}{Lemma}
\usepackage{cleveref}
\crefname{lemma}{Lemma}{Lemmas}
\title{Characteristic Operators and Spectral Properties of Periodic Evolutionary Systems}
\author{Bram Lentjes, Babette A. J. de Wolff}
\date{Source: arXiv:2603.21711v1 (CC BY 4.0)}
\begin{document}
\maketitle

\noindent\emph{Source and licence.} Characteristic Operators and Spectral Properties of Periodic Evolutionary Systems. Version arXiv:2603.21711v1 (CC BY 4.0), distributed by arXiv under the Creative Commons Attribution 4.0 International licence, http://creativecommons.org/licenses/by/4.0/, as recorded in the arXiv licence metadata for this exact version and shown on the abstract page https://arxiv.org/abs/2603.21711v1. Full licence notice: \texttt{licenses/arxiv-2603.21711-CC-BY-4.0.txt}.\par\smallskip

\noindent\textit{Editorial scope.} Counted unit: the article's lemma lemma:DAhatcompo (source0.tex lines 1088--1094). The article's complete original proof of it is delivered with it (source0.tex lines 1095--1125, one \texttt{proof} environment of 31 lines). No mathematics is written, repaired or re-derived: the statement and the proof are the article's own lines, in the article's own notation, and no retained line is substituted or re-flowed.  Retained uncounted as the source-local context the proof needs: lines 756--765; lines 1020--1041; lines 1045--1051; lines 1054--1056; lines 1081--1086.  Nothing was omitted from any retained span, so the package carries no omission record.  No retained line contains a citation, so the wrapper carries no bibliography and no bibliography entry.  No retained line contains a figure environment, an image-inclusion command or drawing code, so no image is delivered and the package declares no asset dependency.  Editorial formatting only, in addition: an AMS \texttt{amsart} preamble that restates the article's own declarations and macros used by the retained lines, namely the environments \texttt{lemma} on separate counters, together with the standard packages those lines need.  The retained environments print in this document as Lemma 1 in order of appearance, whereas the article prints the same items with the numbering of its own full document.  The complete native source of the article as distributed in the arXiv e-print for this exact version is bundled unchanged as \texttt{sources/196969/source0.tex}.
\begin{lemma} \label{lemma:Q(z)closable}
The closable linear operator $Q(z) : \mathcal{D}(Q(z)) \subseteq \mathcal{F}_T(\mathbb{R},X) \to \mathcal{F}_T(\mathbb{R},X)$ defined by
\begin{equation} \label{eq:actionQ}
    \mathcal{D}(Q(z)) \coloneqq \mathcal{D}(D), \quad Q(z) \coloneqq I - R(z,D_0)(zI - D), \quad \forall z \in \Omega \cup \{0\},
\end{equation}
is a projection. Moreover, $Q(z) \iota: \mathcal{Y} \to \mathcal{N}(zI-D)$ is a bounded linear operator with bounded (algebraic) inverse, which can be represented by
\begin{equation} \label{eq:Q(z)iota}
    Q(z)\iota = \iota - zR(z,D_0)\iota, \qquad [Q(z)\iota]^{-1} = \iota^{-1} Q(0) |_{\mathcal{N}(zI-D)}, \quad \forall z \in \Omega \cup \{0\}.
\end{equation}
\end{lemma}

\begin{equation} \label{eq:T12action}
 T_1 
    \begin{pmatrix}
        q \\
        \varphi
    \end{pmatrix}
    \coloneqq JD\varphi = \begin{pmatrix}
        MD \varphi \\
        D \varphi
    \end{pmatrix},    
    \quad 
    T_2 
        \begin{pmatrix}
        q \\
        \varphi
    \end{pmatrix}
    \coloneqq
        \begin{pmatrix}
        (K - MD)\varphi \\
        0
    \end{pmatrix}.
\end{equation}

\begin{lemma} \label{lemma:T2onto}
The following statements hold:
\begin{enumerate}
    \item $\mathcal{R}(\Delta(z)) \subseteq \mathcal{R}(MD - K)$ and $\mathcal{R}(\Delta(z)) \oplus \{0\} \subseteq \mathcal{R}(T_2)$ for all $z \in \Omega$.
    \item If $\rho(\Delta) \neq \emptyset$, then $\mathcal{R}(T_2)$ is dense in $\Gamma(M)^c$.
\end{enumerate}
\end{lemma}

\begin{equation} \label{eq:zM_on_Q}
    z M Q(z) \iota = MD Q(z) \iota,
\end{equation}

To continue, we actually need an equality between the linear (sub)spaces $\mathcal{R}(\Delta(z))$ and $ \mathcal{R}(MD - K)$ of $\mathcal{F}_T(\mathbb{R},Y)$ for at least one $z \in \rho(\Delta)$. Therefore, we impose the following hypothesis:
\begin{enumerate}[label=({H}{{\arabic*}})]
\setcounter{enumi}{5}
\item \label{hyp:rangeequal} 
$\rho(\Delta) \neq \emptyset$ and there exists a $z_0 \in \rho(\Delta)$ such that $\mathcal{R}(\Delta(z_0)) = \mathcal{R}(MD - K)$.
\end{enumerate}

\begin{lemma} \label{lemma:DAhatcompo}
We have the direct sum decomposition
\begin{equation} \label{eq:D(Ahat)directsum}
    \mathcal{D}(\hat{\mathcal{A}}) = \mathcal{D}(J\mathcal{A}J^{-1}) \oplus \mathcal{R}(T_2^+),
\end{equation}
where $T_2^+$ is a right inverse of $T_2$.
\end{lemma}

\begin{proof}
According to \ref{hyp:rangeequal}, let $z_0\in \rho(\Delta)$ be such that $\mathcal{R}(\Delta(z_0)) = \mathcal{R}(MD - K)$. Since $\mathcal{R}(\Delta(z_0)) \oplus \{0 \} \subseteq \mathcal{R}(T_2)$ due to \cref{lemma:T2onto}, and
$\mathcal{R}(T_2) \subseteq \mathcal{R}(MD - K) \oplus \{0 \}$ by \eqref{eq:T12action}, it follows that $\mathcal{R}(T_2) = \mathcal{R}(\Delta(z_0)) \oplus \{0\}$. Consider the linear map $T_2^+ : \mathcal{R}(T_2) \to \mathcal{D}(T_2)$ defined by
\begin{equation*}
    T_2^{+}
    \begin{pmatrix}
        q \\
        0
    \end{pmatrix}
    \coloneqq
    \begin{pmatrix}
        -MQ(z_0)\iota\Delta(z_0)^{-1}q \\
        -Q(z_0)\iota\Delta(z_0)^{-1}q
    \end{pmatrix},
\end{equation*}
which is well-defined since $Q(z_0)$ maps into $\mathcal{D}(D)$, recall \cref{lemma:Q(z)closable}. 
A brief calculation using \eqref{eq:zM_on_Q} and \cref{lemma:Q(z)closable} shows that $T_2 T_2^{+} = I_{\mathcal{R}(T_2)}$, and thus $T_2^+$ is a right inverse of $T_2$. We next introduce the projection operator $P \coloneqq T_2^{+}T_2 |_{\mathcal{D}(\hat{\mathcal{A}})} : \mathcal{D}(\hat{\mathcal{A}}) \to \mathcal{D}(\hat{\mathcal{A}})$ defined by
\begin{equation} \label{eq:projectionP}
    P
    \begin{pmatrix}
        M\varphi \\
        \varphi
    \end{pmatrix}
    \coloneqq
    \begin{pmatrix}
        MQ(z_0)\iota\Delta(z_0)^{-1}(MD-K)\varphi \\
        Q(z_0)\iota\Delta(z_0)^{-1}(MD-K)\varphi
    \end{pmatrix}. 
\end{equation}
To determine the range and kernel of this projection, we first observe that $\mathcal{R}(P) = \mathcal{R}(T_2^+)$ since $\mathcal{R}(T_{2})$ equals the range of $T_2$ restricted to $\mathcal{D}(\hat{\mathcal{A}})$. Consider now the linear operator $I-P : \mathcal{D}(\hat{\mathcal{A}}) \to \mathcal{D}(\hat{\mathcal{A}})$ and note directly that $\mathcal{D}(J\mathcal{A}J^{-1}) \subseteq \mathcal{R}(I-P)$ by a simple calculation. To prove the other inclusion, let $(M\psi,\psi) \in \mathcal{R}(I-P)$ be given. Then $\psi = [I - Q(z_0)\iota \Delta(z_0)^{-1} (MD-K)] \varphi$ for some $\varphi \in \mathcal{D}(D)$. Recalling \eqref{eq:zM_on_Q} shows that $(MD-K)\psi = 0$ and thus $(M \psi,\psi) \in \mathcal{D}(J\mathcal{A}J^{-1})$. We conclude that $P$ is a projection with range $\mathcal{R}(T_2^+)$ and kernel $\mathcal{D}(JAJ^{-1})$ so that direct sum decomposition \eqref{eq:D(Ahat)directsum} follows.
\end{proof}

\end{document}
